\documentclass[11pt]{amsart}
\usepackage{amsmath,amssymb,amsthm}
\usepackage[margin=1.15in]{geometry}
\usepackage[hidelinks]{hyperref}

\newtheorem{theorem}{Theorem}[section]
\newtheorem{lemma}[theorem]{Lemma}
\newtheorem{corollary}[theorem]{Corollary}
\theoremstyle{remark}
\newtheorem{remark}[theorem]{Remark}

\newcommand{\arxiv}[1]{\href{https://arxiv.org/abs/#1}{arXiv:\nolinkurl{#1}}}
\newcommand{\N}{\mathcal{N}}
\newcommand{\cl}[2]{N_{#1}[#2]}
\setlength{\emergencystretch}{2.5em}
\makeatletter\g@addto@macro\UrlBreaks{\do-}\makeatother

\begin{document}

\title[Independent face classes for \v{C}ech complexes]{Independent face classes in \v{C}ech complexes of hypercube
graphs}
\author{Yitzchak Shmalo}
\date{26 November 2026}
\hypersetup{pdftitle={Independent face classes in \v{C}ech complexes of hypercube graphs},pdfauthor={Yitzchak Shmalo}}
\address{Einstein Institute of Mathematics, The Hebrew University of Jerusalem, Jerusalem, Israel}
\email{yspennstate@gmail.com}
\subjclass[2020]{55N31, 55U10, 05C12}
\keywords{\v{C}ech complex, hypercube graph, nerve, homology, retraction}

\begin{abstract}
Let $\N(G,r)$ be the \v{C}ech complex of a graph $G$ at scale $r$. For the hypercube graph $Q_n$ and the scales
$r=2p-2$ with $p\ge2$ and $r=2p-3$ with $p\ge3$, at which $Q_p$ is the smallest cube whose complex is not a simplex, we
show that the $2^{n-p}\binom np$ coordinate $p$-faces of $Q_n$ carry independent split copies of the reduced homology of
$\N(Q_p,r)$, with any coefficients. The retraction onto a face is the coordinate clamp of Adams, Shukla and Singh,
corrected by one bit so that it sends every other face into a simplex. Consequently the homology of $\N(Q_n,2p-2)$ has rank at
least $2^{n-p}\binom np$ in degree $2^p-2$, and $H_4(\N(Q_n,3))$ has rank at least $3\cdot2^{n-3}\binom n3$, which equals the
computed values $24$ and $120$ for $n=4,5$.
\end{abstract}

\maketitle

\section{Introduction}

Let $G$ be a finite connected simple graph, and give its geometric realization the shortest path metric in which every
edge is isometric to $[0,1]$. For $r\ge0$ the \v{C}ech complex $\N(G,r)$ has vertex set $V(G)$, and a set of vertices is
a simplex when the closed balls of radius $r/2$ about them have a common point. Since $\N(G,r)=\N(G,\lfloor r\rfloor)$
for all real $r\ge0$~\cite[Section~2]{ASS}, only integer scales matter. The hypercube graph $Q_n$ has the words in
$\{0,1\}^n$ as vertices, with edges between words at Hamming distance one.

Adams, Shukla and Singh~\cite{ASS} studied $\N(Q_n,r)$; among other results they computed small Betti numbers by
computer~\cite[Table~1]{ASS} and showed that clamping the coordinates outside an $m$-dimensional face gives a retraction
of $\N(Q_n,3)$ onto the complex of that face~\cite[Lemma~5.7]{ASS}, from which $H_3$ and $H_4$ of $\N(Q_n,3)$ are
nonzero for all $n\ge4$~\cite[Theorem~5.13]{ASS}. For Vietoris--Rips complexes of hypercube graphs, Adams and
Virk~\cite[Theorem~6.4]{AV} proved that the homology of the smallest cube carrying a nonzero class propagates: if $p$ is
least with $H_q(\mathrm{VR}(Q_p;r))\ne0$, then the rank of $H_q(\mathrm{VR}(Q_n;r))$ is at least
$\sum_{i=p}^n2^{i-p}\binom{i-1}{p-1}$ times that of $H_q(\mathrm{VR}(Q_p;r))$. They asked which propagation results
hold for \v{C}ech complexes of hypercube graphs~\cite[Question~8.8]{AV}. The following is such a result.

For a set $S$ of $p$ coordinates and a word $b\in\{0,1\}^{[n]\setminus S}$, the coordinate $p$-face $F=F(S,b)$ is the
set of words that agree with $b$ off $S$. Restriction to $S$ is a bijection $\pi_F\colon F\to\{0,1\}^p$ that preserves
Hamming distance, so it identifies $\N(F,r)$, the full subcomplex of $\N(Q_n,r)$ on $F$, with $\N(Q_p,r)$; let
$i_F\colon\N(Q_p,r)\to\N(Q_n,r)$ be the resulting inclusion. There are $2^{n-p}\binom np$ coordinate $p$-faces.

\begin{theorem}\label{thm:main}
Let $r=2p-2$ with $p\ge2$, or $r=2p-3$ with $p\ge3$, and let $n\ge p$. For every coordinate $p$-face $F$ of $Q_n$ there
is a simplicial map $R_F\colon\N(Q_n,r)\to\N(Q_p,r)$ with $R_F\circ i_F=\mathrm{id}$ such that, for every other
coordinate $p$-face $F'$, the image of $R_F\circ i_{F'}$ lies in a single simplex. Consequently, for every abelian group
$M$ and every $q$, the map
\[
\sum_F (i_F)_*\colon\ \bigoplus_F\widetilde H_q(\N(Q_p,r);M)\to\widetilde H_q(\N(Q_n,r);M),
\]
the sum over all $2^{n-p}\binom np$ coordinate $p$-faces, is split injective, with left inverse $\bigoplus_F(R_F)_*$.
\end{theorem}

The consequence follows from the first part at once: $(R_F)_*(i_{F'})_*$ is the identity for $F'=F$, and for $F'\ne F$
it factors through the reduced homology of a simplex, which is zero.

\begin{corollary}\label{cor:rank}
For $p\ge2$ and $n\ge p$, the group $H_{2^p-2}(\N(Q_n,2p-2);\mathbb Z)$ contains a direct summand isomorphic to
$\mathbb Z^{2^{n-p}\binom np}$. For $n\ge3$, $H_4(\N(Q_n,3);\mathbb Z)$ contains a direct summand of rank
$3\cdot2^{n-3}\binom n3$.
\end{corollary}

The first statement uses Lemma~\ref{lem:targets}, by which $\N(Q_p,2p-2)$ is the boundary of a simplex with $2^p$
vertices, a sphere of dimension $2^p-2$. The second uses $\N(Q_3,3)\simeq\bigvee_3S^4$~\cite[Section~5]{ASS}, a
computation the authors attribute to Conant and Thistlethwaite~\cite{CT}. At scale three the bound is $24$ for $n=4$
and $120$ for $n=5$, which are the values of $\beta_4$ in~\cite[Table~1]{ASS}. At scales two and four it lies below the
computed values: $6,24$ against $7,31$ for $\beta_2(\N(Q_n,2))$ with $n=3,4$, and $8,40$ against $10,60$ for
$\beta_6(\N(Q_n,4))$ with $n=4,5$. The count $2^{n-p}\binom np$ includes every $p$-face; the sum in~\cite[Theorem~6.4]{AV}
counts the faces whose fixed coordinates above the largest free coordinate are $0$, which for $n=4$, $p=3$ are $7$ of the
$8$. That $\widetilde H_{2^p-2}(\N(Q_n,2p-2))$ is nonzero for $n\ge p$, the even half of~\cite[Question~7.2]{ASS} in this
degree, already follows from the clamp of~\cite[Lemma~5.7]{ASS}. That lemma is stated at scale three, but the clamp is a
simplicial retraction onto the complex of one face at every scale (Remark~\ref{rem:clamp}). What Theorem~\ref{thm:main}
adds is the rank. It was proved in~\cite{OneStep} that the inclusion of $\N(Q_n,r)$ into $\N(Q_n,r+1)$ is
null-homotopic for all $n$ and $r$; that result says where these classes die, at the next scale, but not how many there are.

At odd scales the theorem is informative exactly when $\widetilde H_*(\N(Q_p,2p-3))\ne0$. By Lemma~\ref{lem:targets} the
simplices of $\N(Q_p,2p-3)$ are the vertex sets whose complement contains an edge, so it is the Alexander dual of the
independence complex of $Q_p$. For $p=3$ its homology is $\mathbb Z^3$ in degree $4$~\cite{CT}. For every $p\ge3$ the
group $\widetilde H_q(\N(Q_p,2p-3);\mathbb Z)$ with $q=3\cdot2^{p-2}-2$ has a free direct summand of rank at least
$3$~\cite[Theorem~1.2]{OddScales}, so by Theorem~\ref{thm:main} the group $\widetilde H_q(\N(Q_n,2p-3);\mathbb Z)$ has a
free direct summand of rank at least $3\cdot2^{n-p}\binom np$ for every $n\ge p$, as noted in~\cite[Section~1]{OddScales}.

\begin{remark}
The hypothesis $p\ge3$ at odd scales cannot be dropped. At scale one, $\N(Q_n,1)$ is the graph $Q_n$ itself, and the
boundaries of the six squares of $Q_3$ are dependent in $H_1$, which has rank $12-8+1=5$.
\end{remark}

\section{The complexes of the face}

Write $d$ for Hamming distance and $\cl k c=\{w:d(w,c)\le k\}$. By~\cite[Lemma~2.2]{ASS}, a vertex set of a graph is a
simplex of $\N(G,2k)$ exactly when it lies in some $\cl kc$, and a simplex of $\N(G,2k+1)$ exactly when it lies in
$\cl ku\cup\cl kv$ for some edge $uv$. A vertex map between graphs therefore induces a simplicial map at scale $r$ as
soon as it sends each of these sets into a simplex. For a word $y$ write $\bar y$ for its antipode, and for an edge
$e=uv$ write $d(w,e)=\min(d(w,u),d(w,v))$ and $\bar e=\bar u\bar v$.

\begin{lemma}\label{lem:targets}
Let $p\ge2$. A set of words in $\{0,1\}^p$ is a simplex of $\N(Q_p,2p-2)$ if and only if it misses some word, and a
simplex of $\N(Q_p,2p-3)$ if and only if it misses both ends of some edge. The complexes $\N(Q_{p-1},2p-2)$ and
$\N(Q_{p-1},2p-3)$ are simplices.
\end{lemma}

\begin{proof}
$\cl{p-1}c$ is everything except $\bar c$. If $uv$ is an edge varying coordinate $i$ and $w$ differs from $u$ in $h$
coordinates other than $i$, then $d(w,uv)=h$ and $d(w,\bar u\bar v)=p-1-h$. So the complement of $\cl{p-2}u\cup\cl{p-2}v$
is the edge $\bar u\bar v$. In $Q_{p-1}$ every word lies within $p-1$ of any vertex and within $p-2$ of any edge.
\end{proof}

\section{Proof of Theorem~\ref{thm:main}}

Fix $F=F(S,b)$ and write $\pi=\pi_F$, extended to all of $Q_n$ by restriction to $S$; then $d(\pi x,\pi y)\le d(x,y)$, and
$d(\pi x,\pi y)\le d(x,y)-1$ when $x\in F$ and $y\notin F$. We define $R_F(x)=\pi x$ for $x\in F$ and
$R_F(x)=\kappa(\pi x)$ for $x\notin F$, where $\kappa\colon\{0,1\}^p\to\{0,1\}^p$ changes at most one bit and is chosen
below; $R_F\circ i_F$ is the identity. The image of a set $A$ splits into the inside part $\pi(A\cap F)$ and the outside
part $\kappa\pi(A\setminus F)$, which lies in the image $K$ of $\kappa$. In both subsections below the argument for another face
$F'=F(S',b')$ starts from the same observation: if $S'=S$ then $b'\ne b$ and $F\cap F'=\emptyset$; otherwise, since
$|S'|=|S|$, there is a coordinate $i\in S\setminus S'$, and $\pi(F\cap F')$ lies in the facet $\{y:y_i=b'_i\}$ of $\{0,1\}^p$.

\subsection*{Even scales}
Let $r=2p-2$ with $p\ge2$, so $k=p-1$. Let $\kappa$ fix the words of even weight and flip the first coordinate of the
others, so $K$ is the set of even words. Consider a generator $\cl kc$ and put $a=\pi c$.

If $c\in F$, then $d(\pi x,a)\le k$ for $x\in F$, and for $x\notin F$ we have $d(\pi x,a)\le k-1$ and hence
$d(\kappa\pi x,a)\le k$. The image lies in $\cl ka$, which misses $\bar a$.

If $c\notin F$, the inside part lies within $k-1=p-2$ of $a$ and so misses the ball $\cl1{\bar a}$. This ball contains a
word of odd weight, which the outside part misses as well.

For $F'\ne F$, the image of $F'$ consists of even words and, if $F\cap F'\ne\emptyset$, of words in a facet
$\{y_i=b'_i\}$. The opposite facet is a cube of dimension $p-1\ge1$ and contains a word of odd weight, which is missed.

\subsection*{Odd scales}
Let $r=2p-3$ with $p\ge3$, so $k=p-2$. Choose three coordinates of $\{0,1\}^p$, call them marked, and let $K$ be the set of
words whose three marked bits are equal. Let $\kappa$ fix $K$ and flip the minority marked bit of every other word. We use
three facts about $K$.
\begin{enumerate}
\item Every word has a neighbour outside $K$: flip a marked bit of a word in $K$, or a majority marked bit of a word
outside $K$.
\item For every edge $e$, the set $\cl1u\cup\cl1v$ ($e=uv$) contains both ends of an edge outside $K$. If the ends of $e$
lie outside $K$ this is $e$. Otherwise the two marked bits that $e$ does not change are equal on both ends of $e$
(if $e$ changes an unmarked coordinate, choose any two of the three); flipping one of them in both ends gives a parallel
edge whose ends lie outside $K$.
\item Every facet $\{y_i=c\}$ contains an edge outside $K$. If $i$ is marked, set a second marked bit to $1-c$ and let
the third vary. If $i$ is unmarked, set two marked bits to $0$ and $1$ and let the third vary.
\end{enumerate}
Consider a generator $A=\cl ku\cup\cl kv$, where the edge $e=uv$ changes coordinate $j$.

If $j\in S$ and $u,v\in F$, then $\pi e$ is an edge, the inside part lies within $k$ of $\pi e$, and the outside part lies
within $k-1$ of $\pi e$ before $\kappa$ and within $k$ after. The image lies in $\cl k{\pi u}\cup\cl k{\pi v}$.

If $j\in S$ and $u,v\notin F$, the inside part lies within $k-1=p-3$ of the edge $\pi e$, hence at distance at least $2$
from the opposite edge $\bar u'\bar v'$, where $\pi e=u'v'$, by the identity in the proof of Lemma~\ref{lem:targets}.
So it misses $\cl1{\bar u'}\cup\cl1{\bar v'}$. By fact~2 this set contains an edge outside $K$, and both of its ends are
missed by the image.

If $j\notin S$, then $\pi u=\pi v=z$ and every word of $\pi(A)$ lies within $k$ of $z$, so the image lies within
$k+1=p-1$ of $z$ and misses $\bar z$. The inside part lies within $p-2$ of $z$ and misses the neighbours of $\bar z$.
By fact~1 some neighbour of $\bar z$ lies outside $K$; it is missed by the image, and so is the edge from it to $\bar z$.

In every case $A$ is sent to a simplex, so $R_F$ is simplicial. For $F'\ne F$, the image of $F'$ consists of words of $K$
and, if $F\cap F'\ne\emptyset$, of words in a facet $\{y_i=b'_i\}$. By fact~3 the opposite facet contains an edge outside
$K$, and both of its ends are missed. This completes the proof of Theorem~\ref{thm:main}.

\begin{remark}\label{rem:clamp}
Without the correction, $R_F=\pi$ is the clamp of~\cite[Lemma~5.7]{ASS}. It is a graph map and is simplicial at every
scale, but it maps every face $F'$ parallel to $F$ onto all of $\{0,1\}^p$, which is not a simplex; the one-bit
correction is what makes the copies from different faces independent.
\end{remark}

\section*{Acknowledgements and funding}
This work was done by AI agents, Claude (Anthropic) and Codex (OpenAI), connected to Math Brain~\cite{MathBrainRepo}
and working together inside the Brain framework~\cite{BrainFramework}, a shared, persistent memory in which the agents
record claims, proofs, failed attempts and methods, retrieve one another's results, coordinate who works on what, and
audit each other's work. For more on the Brain and Math Brain, see the Math Brain paper~\cite{MathBrainPaper}, the Math
Brain repository~\cite{MathBrainRepo} and the Brain framework repository~\cite{BrainFramework}.

The research presented in this paper was supported by the European Research Council (ERC) under the European Union's
Horizon Europe research and innovation programme (grant agreement No.~101041711), by the Simons Foundation as part of
the Collaboration on the Mathematical and Scientific Foundations of Deep Learning, by Heights Labs, by Convex Nexus
Capital, by the Israel Science Foundation (grant number 2258/19), and by the Israel Science Foundation (ISF Grant
4101/25).

\section*{Declaration of competing interest}
The author holds equity in Heights Labs and in Convex Nexus Capital and has a management role at Convex Nexus Capital.
These entities are included in the funding declaration above.

\section*{Data availability}
No data were generated or analysed in this work. A Python script that checks, for every coordinate $p$-face of $Q_n$
in fourteen cases $(p,n)$ with $n\le6$, eight at even scale and six at odd scale, that $R_F$ is simplicial on the
generators of~\cite[Lemma~2.2]{ASS} and that every other face is sent into one simplex, and that the uncorrected clamp
fails the second test, is \texttt{code/check\_face\_retractions.py} at
\url{https://github.com/yspennstate/cech-hypercube-face-classes/tree/v1.3}.

\section*{Declaration of generative AI and AI-assisted technologies}
This paper is part of a research project on AI proving problems. Codex agents (OpenAI) produced the proof, and Claude
agents (Anthropic) checked the proof and the references and wrote the manuscript. The question is due to Adams and
Virk~\cite{AV}, the face retraction at scale three is due to Adams, Shukla and Singh~\cite{ASS}, and the computation
of $\N(Q_3,3)$ is due to Conant and Thistlethwaite~\cite{CT}. The author directed the work and is responsible for the
content of the paper and for decisions concerning its dissemination.

\begin{thebibliography}{9}
\bibitem{ASS} H.~Adams, S.~Shukla and A.~Singh, \emph{\v{C}ech complexes of hypercube graphs}, Homology Homotopy Appl.
\textbf{27} (2025), no.~1, 83--105, \href{https://doi.org/10.4310/HHA.2025.v27.n1.a6}{doi:10.4310/HHA.2025.v27.n1.a6};
\arxiv{2212.05871}. Numbered results are cited as in the arXiv version.
\bibitem{AV} H.~Adams and \v{Z}.~Virk, \emph{Lower bounds on the homology of Vietoris--Rips complexes of hypercube
graphs}, Bull. Malays. Math. Sci. Soc. \textbf{47} (2024), no.~3, Paper No.~72,
\href{https://doi.org/10.1007/s40840-024-01663-x}{doi:10.1007/s40840-024-01663-x}; \arxiv{2309.06222}.
\bibitem{CT} J.~Conant and O.~Thistlethwaite, \emph{Boolean formulae, hypergraphs and combinatorial topology}, Topology
Appl. \textbf{157} (2010), no.~16, 2449--2461.
\bibitem{OddScales} Y.~Shmalo, \emph{Collapsibility numbers and high-degree homology of \v{C}ech complexes of hypercube
graphs}, 2026, \href{https://doi.org/10.5281/zenodo.24569831}{doi:10.5281/zenodo.24569831}.
\bibitem{OneStep} Y.~Shmalo, \emph{A one-step nullhomotopy for \v{C}ech complexes of hypercube graphs}, 2026,
\href{https://doi.org/10.5281/zenodo.24548858}{doi:10.5281/zenodo.24548858}.
\bibitem{MathBrainRepo} Y.~Shmalo, \emph{Math Brain: public snapshot, code and tests}, 2026,
\url{https://github.com/yspennstate/math-brain-public/tree/5427734deee9be643d6dfbdd5e7d4baeeefe4a4a}.
\bibitem{MathBrainPaper} Y.~Shmalo, \emph{Math Brain: shared, auditable memory for multi-agent mathematical
research}, 2026, \href{https://doi.org/10.5281/zenodo.24150837}{doi:10.5281/zenodo.24150837},
\url{https://github.com/yspennstate/math-brain-public/tree/5427734deee9be643d6dfbdd5e7d4baeeefe4a4a/papers/math-brain-2026-10-24}.
\bibitem{BrainFramework} Y.~Shmalo, \emph{Brain Framework}, public software repository, 2026,
\url{https://github.com/yspennstate/brain-framework-public/tree/46b935d26d5a2893f7790b60466387dd8da47ef5}.
\end{thebibliography}

\end{document}
